\subsection{Conclusions}
\label{sec:conclusions}

This thesis investigated whether modern set-based neural networks can separate overlapping electron showers in the LDMX ECal and whether heterogeneous
trigger-pad information adds useful context that these models can extract. Transformer and GravNet models
were trained for fixed \(2e\) and \(3e\) hit assignment, followed by an
experimental TRACE model that predicts slot validity, hit contributors and
energy fractions jointly. The common sharded dataset, deterministic event
splits and held-out test samples made it possible to compare the implemented
models and analyse their failure modes consistently.

Both baseline architectures learned substantial shower separation. The final
Transformer reached permutation-aligned pooled hit accuracies of
\(84.20\,\%\) for \(2e\) events and \(74.84\,\%\) for \(3e\) events.
Energy-weighted accuracy was higher in both samples, which shows that errors
preferentially affect lower-energy hits. The poorer \(3e\) result is
consistent with the smaller closest-pair separation in that sample and with
the observed dependence of accuracy on shower overlap. GravNet performed
similarly, while paired fixed-checkpoint comparisons gave the Transformer a
small advantage of approximately one percentage point or less. The Transformer was
therefore the pragmatic reference architecture in this work, although the
single-seed and implementation-specific comparison does not establish a
universal ranking.

TRACE retained nearly the same aligned dominant-origin performance while supporting a richer reconstruction output with both hit-level and event-level metrics. It achieved
\(99.983\,\%\) electron-count accuracy in the restricted \(2e\)-versus-\(3e\)
setting, while exact contributor sets, mixed-hit identification and energy
fractions remained more difficult. This demonstrates the flexibility of a
particle-flow-inspired set prediction, but also shows that mixed deposits are
the central unresolved part of the problem. Removing \gls{tpad} tokens caused
a small but statistically resolved change in fixed-count hit assignment and a
larger count change in an earlier TRACE study. Since both tests
removed a complete modality only at inference time, they demonstrate
sensitivity rather than the causal value of the detector information.

An additional contribution is the \texttt{ml\_ldmx} software platform
developed for the study. It connects ROOT-data extraction and tensorisation,
sharded datasets, variable-size model inputs, checkpointed training and
reproducible evaluation in one documented workflow. The project is intended
to be reusable by later students and researchers rather than to support only
the models reported here. Taken together, the results establish that
set-based learning is a feasible route to per-electron reconstruction in
simulated \gls{ldmx} pile-up, while shower overlap, mixed hits and validation beyond
the present simulation domain remain the main barriers to application.

\subsection{Outlook}
\label{sec:outlook}

The first priority is a controlled test of heterogeneous detector value.
ECal-only and ECal-plus-\gls{tpad} models should be trained on identical event
splits across several random seeds and evaluated on the same held-out events.
The study should cover a wider range of electron multiplicities and explicitly
separate events with complete, missing and merged \gls{tpad} candidates. An
alternative trigger-scintillator geometry that provides information in both
transverse coordinates (introducing x-coordinates) could also be tested, since it may offer more
complementary context than the present one-dimensional centroids.

TRACE should next replace or reduce its dependence on
canonical \(y\)-ordering. A permutation-invariant loss with event-wise matching,
such as Hungarian assignment, would align predicted slots to truth during
training rather than only during evaluation. Multi-task loss balancing should
then be studied together with broader electron counts and more varied event
conditions. Additional detector streams or reconstructed object-level inputs
can be introduced through controlled modality comparisons so that their
incremental value is measurable.

Before online reconstruction is considered for production, the models should be tested on the
intended hardware using realistic trigger-level inputs. These tests should
measure accuracy, memory consumption and inference time for models of
different sizes. If a smaller or faster model is required, methods such as
pruning, quantisation and knowledge distillation can then be evaluated.

%More specialised attention methods would become relevant if future inputs contained substantially more detector objects. FlashAttention can reduce the memory required for standard self-attention without changing its result. The events studied here contain only a few hundred tokens, so these methods are not an immediate priority.

After robustness and latency have been established, the reconstruction can be
evaluated with trigger-cell inputs and downstream event-level physics metrics.
The documented data, training and evaluation components in \texttt{ml\_ldmx}
provide a handover point for these studies and for other related machine-learning
reconstruction tasks within \gls{ldmx}.